\documentclass[letterpaper]{article}

\usepackage[margin=1.5cm]{geometry}
\usepackage{amsmath,amsfonts}
\usepackage{url}
\usepackage{bm}

\newcommand{\Real}{\mathbb{R}}
\newcommand{\tr}{^\top}
\newcommand{\opt}{^\star}

\setlength{\parskip}{3mm plus1mm minus1mm}
\setlength{\parindent}{0mm}

\title{Problem Set 1}
\author{Foundations of Machine Learning}
\date{}

\begin{document}
\maketitle

\begin{itemize}
\item \textbf{Submission}: Turn in as a \texttt{PDF} and a \textbf{separate source code file} on \url{http://mycourses.unh.edu}. Scanned hand-written notes are encouraged.
\item \textbf{Questions}: See the class website for a link to office hours
\end{itemize}

\textbf{Relevant reading}: IRA 1, 2, 6, 8; MML 5.1; CO 1.1, 3.1.1.  

\section{Recognizing functions}
For each of the following functions $f\colon \Real \to \Real$, state whether it is \emph{linear}, \emph{affine}, \emph{quadratic}, or none of these. Justify each answer in one sentence.
\begin{enumerate}
\item $f(x) = 3x$
\item $f(x) = 3x - 2$
\item $f(x) = x(x-2)$
\item $f(x) = (x+1)^2 - x^2$
\item $f(x) = |x|$
\item $f(x) = 2^x$
\end{enumerate}

\section{Derivatives warm-up}
Compute the first derivative $f'(x)$ and the second derivative $f''(x)$ of each function. Show your work; no credit for unsupported answers.
\begin{enumerate}
\item $f(x) = 4x^3 - 2x^2 + 7$
\item $f(x) = e^{-2x}$
\item $f(x) = \log(1 + x^2)$
\item $f(x) = x e^{x}$
\item $f(x) = \frac{1}{1 + e^{-x}}$ \quad (the sigmoid function; express $f'$ in terms of $f$ if you can)
\item $f(x) = (3x - 1)^4$
\end{enumerate}

\section{Quadratic functions}
Consider the general quadratic function $f(x) = a x^2 + b x + c$ with $a \neq 0$.
\begin{enumerate}
\item Derive the location $x\opt$ of the stationary point of $f$ (the point where $f'(x\opt) = 0$).
\item Under what condition on $a$ is $x\opt$ a minimizer? A maximizer?
\item Write $f(x) = 2x^2 - 8x + 3$ in the form $f(x) = a(x - x\opt)^2 + d$ and use this form to read off its minimizer and minimal value.
\end{enumerate}

\section{Taylor series by hand}
Recall that the first-order and second-order Taylor approximations of $f$ around a point $x_0$ are
\[
f_1(x) = f(x_0) + f'(x_0)(x - x_0), \qquad
f_2(x) = f(x_0) + f'(x_0)(x - x_0) + \tfrac{1}{2} f''(x_0)(x - x_0)^2 .
\]
\begin{enumerate}
\item Compute $f_1$ and $f_2$ for $f(x) = e^x$ around $x_0 = 0$.
\item Compute $f_1$ and $f_2$ for $f(x) = \log(x)$ around $x_0 = 1$.
% \item For each function above, evaluate $f$, $f_1$, and $f_2$ at $x = 0.5$ and at $x = 2$, and report the approximation errors $|f(x) - f_1(x)|$ and $|f(x) - f_2(x)|$. Which approximation is better, and where does the approximation degrade?
\item For which functions $f$ does $f_2 = f$ exactly for every choice of $x_0$? Explain.
\end{enumerate}

\section{Plotting functions \textbf{[code]}}
\begin{enumerate}
\item Plot $f(x) = x^4 - 3x^2 + x$ on the interval $[-2.5, 2.5]$.
\item On the same axes (or a second panel), plot its derivative $f'(x)$. Mark the points where $f'(x) = 0$ and describe, in a sentence or two, how the sign of $f'$ relates to whether $f$ is increasing or decreasing.
\item Plot the sigmoid function $f(x) = 1/(1 + e^{-x})$ on $[-6, 6]$ together with its derivative. Where is the derivative largest?
\end{enumerate}

\section{Visualizing Taylor approximations \textbf{[code]}}
Let $f(x) = \log(1 + e^{x})$.
\begin{enumerate}
\item Derive $f'(x)$ and $f''(x)$ by hand and include the derivation in your writeup.
\item Write a function that, given $x_0$, returns the first-order Taylor approximation $f_1$ and the second-order Taylor approximation $f_2$ of $f$ around $x_0$.
\item Plot $f$, $f_1$, and $f_2$ on $[-4, 4]$ for $x_0 = 0$, and again for $x_0 = 2$, clearly labeling the curves (use a legend).
\item In 2--3 sentences: near $x_0$, which approximation is more accurate? What happens far from $x_0$?
\end{enumerate}

\section{Finding and classifying stationary points}
For each function, find all points where $f'(x) = 0$ and classify each as a local minimum, local maximum, or neither, using the second derivative (or another argument when the second derivative is inconclusive).
\begin{enumerate}
\item $f(x) = x^2 - 6x + 1$
\item $f(x) = x^3 - 3x$
\item $f(x) = x^3$
\item $f(x) = x^4$
% \item $f(x) = x + \frac{4}{x}$ on $x > 0$
\end{enumerate}

% \section{Necessary vs.\ sufficient conditions}
% Answer each question and give a concrete one-dimensional example (a formula and, if helpful, a sketch).
% \begin{enumerate}
% \item Give a function $f$ and a point $x_0$ with $f'(x_0) = 0$ where $x_0$ is \emph{not} a local minimum or maximum.
% \item Give a function on the interval $[0, 1]$ whose minimizer $x\opt$ satisfies $f'(x\opt) \neq 0$. Why does this not contradict the first-order necessary condition?
% \item Give a function with a local minimum that is not a global minimum.
% \item True or false: if $f''(x) > 0$ for all $x \in \Real$, then $f$ has a global minimizer. Prove it or give a counterexample.
% \end{enumerate}

\section{Convexity}
Recall that a twice-differentiable function $f\colon \Real \to \Real$ is convex if and only if $f''(x) \geq 0$ for all $x$.
\begin{enumerate}
\item Determine which of the following functions are convex on the given domain:
  \begin{enumerate}
  \item $f(x) = e^{3x}$ on $\Real$;
  \item $f(x) = \log(x)$ on $x > 0$;
  % \item $f(x) = x^4 - x^2$ on $\Real$; 
  \item $f(x) = \max\{x, 0\}$ on $\Real$ (careful: not differentiable at $0$; argue from the definition or a picture).
  \end{enumerate}
\item Show that if $f$ and $g$ are convex and \emph{twice-differentiable}, then $f + g$ is convex, and $\alpha f$ is convex and for any $\alpha \geq 0$.
% \item Show that $f(x) = (x - 3)^2 + e^{x}$ is convex and explain why any point with $f'(x) = 0$ must be its global minimizer.
\item Give an example of two convex functions whose \emph{elementwise product} is not convex.
\end{enumerate}

% \section{Gradient descent by hand}
% Consider $f(x) = x^2 - 4x + 5$ and the gradient descent iteration $x_{k+1} = x_k - \alpha f'(x_k)$ starting from $x_0 = 0$.
% \begin{enumerate}
% \item Compute $x_1, x_2, x_3$ by hand for step size $\alpha = 0.1$. Report $f(x_k)$ at each iterate.
% \item Repeat for $\alpha = 1$. What happens, and why?
% \item Repeat for $\alpha = 1.5$. What happens, and why?
% \item For this particular $f$, derive the range of step sizes $\alpha$ for which the iteration converges to the minimizer from any starting point. (Hint: write $x_{k+1} - x\opt$ as a multiple of $x_k - x\opt$.)
% \end{enumerate}

% \section{Implementing gradient descent \textbf{[code]}}
% Implement gradient descent for functions of one variable:
% \begin{verbatim}
% gradient_descent(fprime, x0, alpha, iterations) -> vector of iterates
% \end{verbatim}
% where \texttt{fprime} is the derivative of the objective, \texttt{x0} is the starting point, \texttt{alpha} is the step size, and \texttt{iterations} is the number of steps. The function should return all iterates $x_0, x_1, \dots$.
% \begin{enumerate}
% \item Verify your implementation on $f(x) = x^2 - 4x + 5$ from with $\alpha = 0.1$ and $x_0 = 0$, your first three iterates must match your hand computation.
% \item Run your implementation on $f(x) = x^4 - 3x^2 + x$ (from Problem 5) with $\alpha = 0.01$, 200 iterations, from the two starting points $x_0 = -2$ and $x_0 = 2$. Report the final iterate and the final objective value for each. Do the two runs find the same point? Relate what you see to your plot from Problem 5.
% \item For $x_0 = 2$, plot $f(x_k)$ as a function of the iteration $k$ for step sizes $\alpha \in \{0.001, 0.01, 0.1\}$ on the same axes. Describe the trade-off you observe.
% \item Plot the function $f$ and overlay the iterates $(x_k, f(x_k))$ for $\alpha = 0.01$, $x_0 = 2$, so that the path of the algorithm is visible.
% \end{enumerate}

% \section{Gradient descent meets Taylor \textbf{[code]}}
% This problem connects the two halves of the assignment. Let $f(x) = \log(1 + e^{x}) + \tfrac{1}{2}(x-1)^2$.
% \begin{enumerate}
% \item Argue (by hand) that $f$ is convex.
% \item Run your gradient descent from Problem 11 with $\alpha = 0.5$ and $x_0 = -3$ until the iterates stop changing noticeably; report the (approximate) minimizer $x\opt$.
% \item Produce a single plot showing $f$, its first-order Taylor approximation around $x_0 = -3$, and its second-order Taylor approximation around $x_0 = -3$, with the gradient descent iterates marked on the graph of $f$.
% \item In 2--3 sentences: the gradient descent step moves in the direction in which the \emph{first-order} Taylor approximation decreases. Explain why the step size must then be kept small.
% \end{enumerate}

\section{Time spent}
Report approximately how many hours you spent completing this assignment (including the reading).

\end{document}

%%% Local Variables:
%%% mode: LaTeX
%%% TeX-master: t
%%% End:
